\chapter{CONCLUSION ET PERSPECTIVES}
\label{ch:conclusion}

%------------------------------------------------------------------------------
\section{Bilan du projet}
%------------------------------------------------------------------------------

Ce Projet de Fin d'Études, conduit chez \textbf{SEVENAPP} en mission \textbf{consultant Data} chez \textbf{EQDOM}, a permis de concevoir et d'implémenter progressivement une \textbf{Cloud Data Platform} alignée sur un cahier des charges académique exigeant : architecture Lakehouse médaillon, Kubernetes operators, Infrastructure as Code, sécurité, observabilité et optimisation des coûts.

Les apports principaux :
\begin{itemize}
    \item une \textbf{landing zone AWS/EKS} modulaire (Terraform, IRSA, KMS, deux node groups) ;
    \item une \textbf{chaîne data bout en bout} documentée et préparée (CDC, Iceberg, dbt, Dremio, Power BI) ;
    \item une \textbf{industrialisation} par scripts, CI/CD et ADR ;
    \item une \textbf{traçabilité} explicite pour la soutenance (\texttt{SPEC\_ALIGNMENT.md}) ;
    \item une \textbf{couche d'observabilité} couvrant monitoring, logging et dashboarding technique.
\end{itemize}

%------------------------------------------------------------------------------
\section{Limites}
%------------------------------------------------------------------------------

\begin{itemize}
    \item Environnement \texttt{dev} unique ; haute disponibilité production (NAT multi-AZ, backups CNPG/Barman) non déployée ;
    \item Ansible et ArgoCD restent des extensions ou scaffolds ;
    \item validation E2E et jeux de données réels EQDOM hors périmètre strict du MVP académique ;
    \item TLS production (Let's Encrypt / ACME) à durcir au-delà du MVP cert-manager.
\end{itemize}

%------------------------------------------------------------------------------
\section{Perspectives}
%------------------------------------------------------------------------------

\begin{itemize}
    \item Passage en \texttt{prod} : multi-NAT, sauvegardes, politiques de rétention Iceberg ;
    \item Activation \texttt{lambda\_scaling\_enabled} pour optimisation des coûts ;
    \item Intégration Power BI, dashboards métier et catalogues de données EQDOM ;
    \item GitOps ArgoCD et dépôt d'applications dédié ;
    \item Renforcement qualité : tests automatisés pipeline, Data Quality (Great Expectations), alerting fonctionnel sur fraîcheur et complétude des données.
\end{itemize}

%------------------------------------------------------------------------------
\section{Apports personnels}
%------------------------------------------------------------------------------

Ce stage a consolidé les compétences en \textbf{data engineering cloud-native} : Terraform, Kubernetes operators, streaming CDC et architecture médaillon — compétences directement transférables aux projets de transformation data du secteur financier.

\vspace{1cm}
\hfill \NomEtudiant

\clearpage
